% The learners every model is given: the same ladder of runs on every task,
% so each paper reads it from here under its own section heading.
A run is one learner, trained on the train rows, tuned on the valid rows and
scored once on the test rows. Every model keeps the same ladder of runs, so
each step up it says what one more thing is worth. Table~\ref{tab:runs} lists
this model's.

\paragraph{The floor.}
\run{constant} says one number for every row: the mean label of the train
rows, or their median where the score is an absolute error. A model that does
not beat it has learned nothing.

\paragraph{Trees.}
\run{gbm-context} is gradient-boosted trees (LightGBM~\cite{lightgbm}) on the
context numbers alone, with no view of the board. It measures how much of the
answer is already in the numbers beside the board. \run{gbm-hand} adds the
hand features and is the \emph{reference}: every other run is compared with
it on the same test rows. Both grow up to \V{method.gbm.rounds} trees of up
to \V{method.gbm.leaves} leaves at a learning rate of \V{method.gbm.rate},
each tree on a random \V{method.gbm.rows.percent}\% of the rows and
\V{method.gbm.features.percent}\% of the features, and stop when the valid
rows have not improved for \V{method.gbm.patience} rounds.

\paragraph{Networks.}
A network is a body, which turns the tokens, the marks and the context numbers
into one vector, and a head of two layers, which turns that vector into the
answer. Two bodies ship, so that a leaderboard can say what the architecture
is worth on a task.

The \emph{square transformer}~\cite{attention} treats each square as a token:
an embedding of its piece, plus a learned embedding of the square, plus a
linear map of the square's marks. Every square attends to every other from the
first layer, so a bishop on a1 and a king on h8 meet at once. Attention is
told about the board's geometry by a learned bias for each offset between two
squares, one per head~\cite{relative}: files, ranks, diagonals and knight hops
are all offsets, so a head can learn to look along them. The context numbers
enter as one more token, and the body's output is that token beside the mean
of the 64 square tokens.

The \emph{convolutional tower} is residual blocks of three-by-three
convolutions over the board, the kind of network the first AlphaZero-style
engines used~\cite{alphazero}. The context numbers are added to every square
after the first layer, and the output is the mean and the maximum over the
squares. It sees further only as it gets deeper. It is cheap and trains on a
CPU, and it is here as the honest alternative: where it matches the
transformer, the cheaper model wins.

Sizes are settings of a run, not code. \run{conv-small} and
\run{transformer-small} are sized to train on a CPU. \run{transformer-medium}
is sized for a GPU, and \run{transformer-medium-hand} is the same network
reading the hand features beside the board. Table~\ref{tab:settings} in the
appendix has every size and schedule.

\paragraph{One training loop.}
Every network is trained the same way: AdamW~\cite{adamw}, a learning rate
that rises over the first \V{networks.warmup.percent}\% of the steps and then
follows a cosine down to zero, gradients clipped, mixed precision on a GPU.
The loss follows the score: squared error for a quantity ranked by RMSE, a
smoothed absolute error for one ranked by MAE, the log loss for a chance. A
quantity's label is standardised by the train rows. After each epoch the
network is scored on the valid rows, the weights of the best epoch are the
ones kept, and training stops early once \V{networks.patience} epochs pass
without a better one.

\begin{table}[htbp]\centering\footnotesize
\block{runs}
\caption{The runs, best first. \emph{Kept} is what the valid rows chose: the
number of trees, or the epoch whose weights were kept. \emph{Time} is the
whole run, training and scoring.}
\label{tab:runs}
\end{table}
